\documentclass[10pt,letterpaper]{article}

\usepackage{geometry}
\geometry{margin=1in}
\usepackage{enumitem}
\usepackage{setspace}
\usepackage{titlesec}
\usepackage[colorlinks=true, linkcolor=black, citecolor=blue, urlcolor=blue]{hyperref}
\usepackage[none]{hyphenat}
\sloppy
\usepackage[bottom,hang,flushmargin]{footmisc}
\setlength{\footnotemargin}{1em}

\makeatletter
\renewcommand\footnoterule{%
  \kern -3pt
  \hrule width 0.4\columnwidth height 0.4pt
  \kern 1em
}
\makeatother

\usepackage{fontspec}
\setmainfont[
    Path=fonts/,
    UprightFont = *-Regular,
    BoldFont = *-Bold,
    ItalicFont = *-Italic,
    BoldItalicFont = *-BoldItalic
]{NoticiaText}

\renewcommand{\baselinestretch}{1.15}

\makeatletter
\renewcommand\normalsize{%
   \@setfontsize\normalsize{10pt}{14pt}
}
\makeatother

\title{
  The Adequacy of the Bangladesh Copyright Act, 2023 in Protecting Digital Works and Addressing Piracy
}

\author{Md. Shariar Arfin}
\date{}

\begin{document}
\maketitle

\vspace{1em}

\section{Introduction}
Over the past decade, Bangladesh, like many other developing countries, has experienced a major transformation in the way creative works are produced, distributed, and consumed. Digital technologies have opened new ways for writers, artists, and other creators to reach audiences, but at the same time, they have created difficulties for the law in keeping pace with these changes. To address this situation, the Government of Bangladesh enacted the Copyright Act, 2023 to address the challenges posed by digital works and piracy. This paper examines the adequacy of the Bangladesh Copyright Act, 2023 in protecting digital works and combating piracy.

% A nation's commitment to protecting digital works is fundamentally assessed through its alignment with international treaty standards, which provide the jurisdictional and substantive framework for cross-border protection.

% 2. Literature Review and Theoretical Framework
%   • Evolution of Copyright Law in the Digital Age
% \subsection{International Conventions and Treaties on Digital Copyright}
%   • Comparative Perspectives on Digital Copyright Legislation
%   • Impact of Piracy on Creative and Digital Industries

\section{Key Provisions Related to Digital Works}
Digital work based on information technology- means a creative work directly or indirectly created and used on such devices, such as computers, mobile phones or any digital device, etc., by processing information and data for the purpose of obtaining a specific result, including information, source code, tables, charts, graphs, sounds, images, moving images, designs, text, instructions, signals and instructions for using such works.\footnote{CA 2023, s 2(10)}

The author or proprietor of an information technology-based digital work is vested with a range of exclusive rights. These includes any action application to literary, dramatic, musical, artistic, cinematic, and sound works, which includes reproduction, performance, publication, translation, adaptation, and broadcasting across various forms of expression. Furthermore, the rights holder retains the exclusive authority to sell, rent copies, provide service, license regardless of whether similar copies were previously sold or rented.\footnote{CA 2023, s 2(7)(b)}


\section{Addressing Digital Piracy and Penalties}
If any person other than the legitimate owner of the copyright or related rights, engages in the reproduction or creation of an entire work creator is called copyright infringement which is commonly known as piracy. This act involves using any electromagnetic, digital, or other mechanical means to copy the entirety or even a part of the work, including making, importing, or communicating a complete or partial duplicate without proper authorization.\footnote{CA 2023, s 69}

If it appears that the copyright or other rights of the author of a work or its legal license have been infringed by a network service provider or any other third party, the author of the work or its legal license or the network service provider or the third party shall, subject to written objection to the author of the work or the person or institution holding the right to do so or the network service provider or the third party, remove all copies of the objected work from any medium under its control as soon as possible and notify the complainant in writing; otherwise, the network or institution, person or third party providing the service shall be liable for the crime of copyright infringement. In the case of commercial use of any work, any institution shall have to obtain a license from the Bangladesh Copyright Office.\footnote{CA 2023, s 82}

If any person infringes any information technology-based digital work and uses, publishes, sells or distributes a copy thereof through any medium, it shall be an offence and he shall be punished with imprisonment for a term not exceeding 4 (four) years and a fine not exceeding 4 (four) lakh taka. Provided that if it is proved to the satisfaction of the court that the information technology-based digital work was not infringed for the purpose of making profit in the course of business or commercial activities, then the person guilty of the said offence shall be punished with imprisonment for a term not less than 3 (three) months or a fine not less than 25 (twenty-five) thousand taka.\footnote{CA 2023, s 100}

If any person convicted under this Act is again guilty of a similar offence, he shall be punished for the second and every subsequent offence with imprisonment not exceeding 5 (five) years or with a fine not exceeding 5 (five) lakh taka.\footnote{CA 2023, s 106}

\section{Enforcement Challenges in Protecting Digital Works}

Despite the robust provisions introduced by the Copyright Act, 2023, the practical enforcement of digital copyright protection in Bangladesh faces numerous persistent obstacles:

\begin{enumerate}
    \item \textbf{Weak Enforcement Mechanisms and Institutional Capacity}
    \begin{itemize}
        \item \textit{Lack of Specialized Expertise:} 
        Law enforcement agencies and prosecutors frequently lack the technical training and specialized resources necessary to effectively investigate and prosecute complex online piracy. This deficiency specifically impacts the ability to track encrypted transfers, identify sources behind anonymized servers, perform forensic analysis of software code, and analyze cloud infrastructure and network tracking data.\footnote{Nafew Sajed Joy, `Why copyright law reform is crucial for Bangladesh' (The Business Standard, 21 November 2024) <https://www.tbsnews.net/thoughts/why-copyright-law-reform-crucial-bangladesh-999156> accessed 21 October 2025.} This functional gap allows digital infringers to operate with relative impunity, bypassing conventional investigative methods.

        \item \textit{Task Force Implementation:} Although the law provides for a national Copyright Task Force, its effective formation and operationalization—requiring coordination across government, law enforcement, and industry stakeholders—remains a work in progress.
    \end{itemize}

    \item \textbf{Difficulty in Tackling Borderless Piracy}
    \begin{itemize}
        \item \textit{Jurisdictional Barriers and Difficulty in Tracking:} The enforcement of digital copyright is critically challenged by the transnational nature of online piracy. The distribution of infringing content via global networks, often utilizing overseas servers and anonymized domain registration, creates significant jurisdictional barriers.\footnote{Leppard Law, `How Digital Piracy Challenges Copyright Enforcement Across Borders' accessed 21 October 2025.} This structure fundamentally impedes Bangladeshi authorities from asserting domestic legal authority, securing essential international cooperation, and successfully issuing and executing cross-border takedown orders and prosecuting responsible parties.
   
        \item \textit{Digital Rights Management (DRM) Circumvention:} Persistent efforts by piracy groups to bypass DRM and other technological protection measures (TPMs) undermine creators’ ability to protect digital works and escalate the costs of enforcement.
    \end{itemize}

    \item \textbf{Limited Public Awareness and Consumer Behavior}
    \begin{itemize}
        \item \textit{Piracy Mindset:} There is a widespread lack of public understanding of copyright as a fundamental property right. Coupled with cultural acceptance of not paying for digital content, the mass consumption of pirated software, music, and media remains the norm.\footnote{Rahman MA, `Designing Copyright Laws to Combat Digital Piracy and Effectively Balance Proprietary and Public Interests in Bangladesh'  accessed 21 October 2025}
        \item \textit{Non-Commercial Loophole:} Although the 2023 Act seeks to address previous loopholes that allowed lighter penalties for ``non-commercial'' infringements, this legacy distinction continues to blur the perceived seriousness of piracy among the public.
    \end{itemize}

    \item \textbf{Challenges for Digital Platforms (ISPs and OTTs)}
    \begin{itemize}
        \item \textit{ISP Liability and Takedowns:} Internet and digital service providers are now explicitly required to remove infringing material once notified, but the volume and speed of new infringements demand constant surveillance and rapid response, making enforcement costly and complex.
        \item \textit{Evolving Distribution Methods:} Piracy networks frequently shift tactics, distributing illicit content through OTT platforms, messaging apps (e.g., Telegram channels), and unauthorized APK portals, requiring regulators and rights holders to continually innovate their monitoring and enforcement strategies.\footnote{The Financial Express `OTT platforms urge BTRC action on illegal streaming sites'  accessed 21 October 2025}
    \end{itemize}

   \item \textbf{Emerging Challenges and Deficiency of Deterrent Measures}
   \begin{itemize}
      \item \textit{Inadequate Remedies and Deterrence:}
      Despite recent legislative reforms, the efficacy of the enforcement regime is significantly hampered by historically lenient penalty structures and ineffective remedial measures against proven copyright infringement. This deficiency in the legal framework compromises deterrence, failing to impose costs proportional to the scale of digital piracy and consequently encouraging repeated violations that undermine the commercial viability and market stability of protected works.

      \item \textit{Emerging Technologies and Challenges}
      The rapid and continuous evolution of digital platforms presents a highly dynamic challenge to effective copyright administration. While platforms such as streaming services, social media, and peer-to-peer (P2P) networks facilitate legitimate distribution, they are simultaneously exploited to expedite the unauthorized, mass distribution of copyrighted material without adequate technical or regulatory safeguards. Furthermore, the advent of technologies like Generative Artificial Intelligence (AI) introduces profound and novel complexities. AI tools challenge conventional concepts of authorship and originality when training on copyrighted content, and they enable the creation of high-quality, scalable derivative works or deepfakes that exponentially increase the potential for automated infringement, necessitating proactive legal and technical adaptation.
   \end{itemize}
\end{enumerate}


\section{Proposed Reforms to Strengthen Copyright Enforcement}
To transition from mere legislative compliance to robust enforcement, research suggests a multi-pronged strategy focusing on judicial, technical, and cultural reforms:

\begin{itemize}
   \item \textbf{Establishment of Specialized Tribunals:} The primary recommendation is the creation of dedicated Intellectual Property (IP) Tribunals or Courts to ensure expedited, expert, and consistent adjudication of digital infringement cases. This specialization is critical for building a predictable IP enforcement environment.
   
   \item \textbf{Investment in Proactive Mechanisms:} The government, in collaboration with industry, should invest in automated digital monitoring and web-scraping tools capable of continuously tracking the unauthorized distribution of digital works. Such proactive enforcement mechanisms are necessary to stay ahead of evolving piracy technology.
   
   \item \textbf{Systematic Public Education:} Sustained, well-funded public awareness campaigns must be initiated across media and educational institutions to systematically educate the populace about intellectual property ethics and the direct economic benefits of consuming legitimate content.
   
   \item \textbf{Empowerment of the Copyright Office:} The Copyright Office's operational capacity should be expanded to transform it from a purely registration-based body into a proactive enforcement agency. This involves granting the Copyright Office expanded authority to investigate infringement cases, issue administrative takedown notices, and provide comprehensive technical and legal support to creators engaged in anti-piracy efforts.
   
   \item \textbf{Compulsory Technical Training:} Mandatory, continuous training must be provided to law enforcement officials, public prosecutors, and customs authorities focused on digital forensics and cyber investigation techniques to close the existing expertise gap.
   
   \item \textbf{Inter-Agency Coordination and Cooperation:}
   The effective transition from updated legislation to practical digital enforcement necessitates resolving a systemic deficiency in inter-agency coordination. Despite the provision for a Copyright Task Force within the legal framework, its effective operationalization is hampered by a lack of seamless cooperation and resource-sharing among key bodies, including the Copyright Office, law enforcement agencies, and the Bangladesh Telecommunication Regulatory Commission. This fragmentation of effort results in slow response times, disjointed investigations, and a failure to swiftly implement comprehensive site-blocking and content-removal strategies, which is critical for mitigating the large-scale, rapid distribution characteristic of online piracy.
\end{itemize}


\section{Conclusion}
From the above discussion, it can be concluded that the Copyright Act, 2023 has been an important upgrade of Bangladesh law so far as the statutory recognition and punitive aspect of online works are concerned. Its effectiveness is, however, severely limited by a substantial ``gap between policy and practice". Systemic issues That the Act is ineffective is attributed to ongoing systemic problems including in-house technical expertise among enforcement bodies, convoluted transnational jurisdictions and fragmented inter-agency cooperation. In the end, while the law is strong, whether or not it will succeed in creating a stable ecosystem for digital content and in effectively dissuading advanced form of piracy (especially against modern threats as Generative AIs) will depend entirely on an immediate sustained investment into judicial specialisation, technical capability reinforcement, institutional empowerment in action.
\end{document}
